\documentclass[11pt]{article}

\usepackage{amsmath}
\usepackage{amssymb}
\usepackage[margin=1in]{geometry}
\usepackage{booktabs}
\usepackage{hyperref}

\title{Photon communication in the Monte Carlo module}
\author{Athena++ Monte Carlo radiation transport}
\date{}

\begin{document}
\maketitle

\begin{abstract}
This note records how photons move between MeshBlocks and between MPI ranks, and how the
transport loop decides that every photon has finished. It is written to be read alongside
\texttt{src/monte\_carlo/mcexchange.cpp}, \texttt{src/monte\_carlo/photon.cpp} and the
transport driver in \texttt{src/monte\_carlo/montecarlo.cpp}. The emphasis is on the
invariants each layer maintains, because the failure mode of this subsystem is not a crash
but the silent loss of photons: every defect found while building it was invisible except
through the conservation identity of Sec.~\ref{sec:conservation}. Measured costs and the
history of how the scheme arrived at its present form are recorded separately, in
\texttt{doc/monte\_carlo/transport\_loop\_performance.md}.
\end{abstract}

\section{The problem}

Photon transport is a sweep over MeshBlocks. A photon is pushed until it reaches an
interaction point, leaves the domain, or leaves its block; in the last case it must be
handed to the neighbouring block, which may live on another rank. Transport is therefore a
sequence of \emph{rounds}, each round pushing every photon that is currently resident and
then moving those that crossed a block face.

Two properties of the physics set the shape of the problem.

First, a photon's block-to-block path is unbounded. In an optically thin region with
periodic sides a nearly horizontal photon crosses block after block without reaching an
interaction point, so the number of rounds is set not by the number of photons but by the
longest single trajectory. On a $4096$-block test problem, $99.6\%$ of rounds transported a
single photon.

Second, the work is not spread evenly. Most photons finish quickly and a few do not, so
after the first few rounds almost every rank is idle and waiting on one straggler.

Together these mean the per-round fixed cost, not the per-photon cost, is what the run is
made of. The scheme described here is built to make an idle round as close to free as
possible.

\section{Layers}

Three layers sit between a photon leaving a block and arriving in the next.

\begin{description}
\item[Block to block, same rank.] A direct write into the target block's receive buffer
during the send sweep. No message, no handshake.

\item[Rank to rank.] Everything leaving this rank for a given peer, from every block, is
collected into one buffer per peer by \texttt{MCRankExchange} and sent as one message set.
A block no longer talks to a neighbouring block across a rank boundary; a rank talks to a
rank.

\item[Termination.] A decision, taken by all ranks, that no photon remains anywhere.
\end{description}

The first two are always in force. The third has two implementations, selected by
\texttt{<montecarlo>/async\_term}, described in Secs.~\ref{sec:sync} and \ref{sec:async}.

\section{Buffer layout and the aggregation}
\label{sec:layout}

A photon's state is stored as a struct of arrays: \texttt{nint} integer properties,
\texttt{nreal} real properties and \texttt{ncplx} complex properties per photon, the last
being the polarization tensor and present only in polarized runs. Packing a photon for
transfer copies these into the three streams of a \texttt{ParticleBuffer}.

The real stream carries \texttt{nreal} values per photon, not the $2\,\texttt{nreal}$ the
underlying \texttt{Particles} class provides for. The second copy, \texttt{rp1}, is the
other register of the dust-particle time integrator. A photon is pushed by
\texttt{PhotonPusher} and never reads it: filling \texttt{rp1} with NaN on every resize left
the \texttt{amr\_transport} results bit-identical, which is what established this rather
than an argument from inspection. Carrying it regardless cost half the real payload of every
photon message and half the real storage of every resident photon, so in a Monte Carlo build
it is neither sized nor sent.

\texttt{MCRankExchange::Stage} appends one block's outgoing photons to the buffers for
their destination rank, recording a header triple
\begin{equation}
  (\,\texttt{lid},\ \texttt{bufid},\ \texttt{npar}\,)
  \label{eq:triple}
\end{equation}
where \texttt{lid} is the local block index on the \emph{receiving} rank and \texttt{bufid}
the boundary slot that block expects them in. The three property streams are appended back
to back in the same block order. A message is therefore self-describing: given the headers,
a receiver can walk the streams by advancing three offsets, each by
\texttt{npar} times the corresponding per-photon property count.

That walk is the one piece of arithmetic both transports share, and it is written once, in
\texttt{MCRankExchange::Deliver}. This is deliberate. Its offsets are trivially wrong in
ways that produce no error, only photons delivered with another photon's properties.

\subsection{Peer set}

The set of ranks this one shares a block boundary with is computed once, after the
neighbour lists are linked. The relation is symmetric --- if a block here neighbours a
block there, that block neighbours this one --- so both sides agree on the peer set without
negotiating it, which is what allows a paired exchange with no discovery step.

The "once" is an assumption that holds only because the module runs as static
post-processing on a mesh that never remeshes. Anything that moves blocks between ranks
invalidates both the peer list and every \texttt{lid} in a staged header.

\section{Synchronous termination}
\label{sec:sync}

The original scheme, retained as \texttt{async\_term = false}.

Each round ends with \texttt{FinishRound}, which performs the rank exchange and then a
blocking \texttt{MPI\_Allreduce} over two counts: photons resident on this rank, and
photons still to be emitted. The run continues while either is nonzero anywhere.

The rank exchange here is a two-phase collective over the peer set. Each peer is first told
how much is coming, as four integers, because it cannot post receives of the right size
otherwise; the payload follows only between peers that have one. The size message goes out
even when nothing follows, so an idle round still costs $2 N_{\rm peer}$ messages on every
rank, plus the reduction.

This is correct and simple, and its cost is a constant per round. Given that the round
count is set by the longest trajectory, a constant per round is precisely the wrong thing
to be paying.

\section{Counter-based asynchronous termination}
\label{sec:async}

The default, \texttt{async\_term = true}. No rank waits on any other. A rank transports
what it holds, posts what is leaving without waiting for it to be taken, and picks up
whatever has arrived.

\subsection{Sending and receiving without a handshake}

The size handshake exists only so a receiver knows how much to expect. That information is
already in the header, so the int stream is sent as a single message that leads with its
own triple count:
\begin{equation}
  [\,n_{\rm trip},\ (\texttt{lid},\texttt{bufid},\texttt{npar})^{n_{\rm trip}},\
  \text{int payload}\,] .
  \label{eq:msg}
\end{equation}
A receiver probes for this one message, reads $n_{\rm trip}$, sums the \texttt{npar} to get
the total photon count, and from that computes the exact lengths of the real and complex
streams waiting behind it. MPI preserves message order between a given pair of ranks on a
tag, so those streams belong to the header just read.

The consequence is the point of the exercise: a rank with nothing to send sends nothing at
all. Silence is the signal, and an idle round costs no messages.

\subsection{The termination question}

What remains is deciding when everyone has finished. Two counters answer it. Let $S$ be the
cumulative number of photons this rank has handed to other ranks, incremented at
\texttt{Stage}, and $R$ the number it has taken from them, incremented on delivery. Summing
over ranks,
\begin{equation}
  \sum_r S_r - \sum_r R_r = \text{photons in flight} \ \geq 0 ,
  \label{eq:inflight}
\end{equation}
so the mesh is finished when no rank holds a photon and the two global sums agree.

Local activity must count three things, not one:
\begin{equation}
  A_r = \sum_{\rm blocks} \big( \texttt{nphot} + \texttt{nphremain} \big)
        + \#\{\text{blocks with } \texttt{has\_incoming\_}\} .
  \label{eq:active}
\end{equation}
The third term is not redundant. A photon that has been delivered into a block's receive
buffer but not yet flushed into the block is neither in flight, by
Eq.~\eqref{eq:inflight}, nor visible in \texttt{nphot}. Omitting it lets a photon vanish
from both sides of the test at once, and the run terminates holding it.

\subsection{Why one reduction is not enough}

The test is taken with a nonblocking reduction of $(A, S, R)$, so a rank keeps working
while it is in flight. A single reduction cannot be trusted, because its inputs are read at
different instants on different ranks: rank $b$ may contribute after receiving a photon
that rank $a$ contributed before sending, and the sum then reports $\sum S = \sum R$ with a
photon still on the wire.

Two consecutive reductions settle it. A rank posts the second only after the first has
completed, so every rank's contribution to the second is taken after every rank had
contributed to the first. If both report $A = 0$ globally and the \emph{identical} pair of
totals $(S, R)$ with $S = R$, then no rank sent anything between them and nothing was in
flight. Comparing the totals across the two is essential and not merely defensive: a rank
that received a photon, finished it and fell idle again would otherwise show $A = 0$ twice
in a row while the counters moved underneath. This is Mattern's four-counter argument in
its usual practical form.

\subsection{Ordering constraints}

Three orderings in the loop are load-bearing, each corresponding to a way the scheme
deadlocks or loses photons if written the obvious way.

\begin{enumerate}
\item \textbf{Completing sends must drain.} A message large enough to go by rendezvous does
not complete until the receiver has matched it. A rank that waits on its own sends without
receiving is waiting on a peer doing the same thing. This is not a rare race: it is the
first round, when every rank has a full complement of photons to hand over and all of them
wait at once.

\item \textbf{The wait belongs before the next reset, not after posting.} The send buffers
cannot be reused until the sends release them, but a rank that posted nothing must still
reach the drain, so the wait sits at the top of the following pass rather than at the
bottom of this one.

\item \textbf{"What landed" must be read from the counter.} Because completing sends also
delivers photons, the number that arrived is the change in $R$ across both calls, not the
return value of the second. Taking only the latter skips the flush, leaves those photons in
receive buffers where Eq.~\eqref{eq:active} cannot see them, and terminates the run holding
them.
\end{enumerate}

\subsection{The idle path}

A rank waiting on someone else's straggler goes round the loop as fast as the probe
returns. Two things keep that cheap. It skips both the transport sweep and the activity
scan of Eq.~\eqref{eq:active}, which are $O(N_{\rm block})$ and provably zero when nothing
is resident and nothing arrived. And after a short spin it yields the core, since holding
one at full tilt costs exactly the cycles the rank still working needs.

\section{The conservation identity}
\label{sec:conservation}

Every photon ends in exactly one terminal state, so at the end of a run
\begin{equation}
  n_{\rm esc} + n_{\rm abs} + n_{\rm des} + n_{\rm rem} = n_{\rm tot} ,
  \label{eq:conservation}
\end{equation}
with $n_{\rm rem}$ counting photons retired by the path cap of
Sec.~\ref{sec:capmove}. All five are reported in the run summary.

This identity is the whole test suite for this subsystem. A photon dropped in transfer, or
delivered twice, or left in a buffer at termination, changes nothing that any physics test
would notice --- the spectrum is still a spectrum, drawn from slightly fewer samples --- and
shows up here and nowhere else. Every bug found while building the scheme was caught by it:
a receive that overwrote instead of appending, losing $64\%$ of photons at $16$ ranks; and a
premature termination that left $75\%$ of them sitting undelivered. It should be checked
whenever this code is touched, at more than one rank count, since the failures are
sensitive to how work happens to be distributed.

\section{Bounding a trajectory}
\label{sec:capmove}

The round count is set by the longest single flight, so a cap on that flight is a cap on
the run. \texttt{<montecarlo>/capmove} bounds the number of pusher steps in one free
flight, counting in a photon property so the count survives both the end of a \texttt{Move}
and a transfer to another block, and resetting at each scattering, since a scattering
begins a new flight.

Two details decide whether it works at all. The status test must accept \texttt{BUFFERED}
as well as \texttt{EVOLVING}: a \texttt{Move} ending \texttt{EVOLVING} is one that reached
an interaction point inside the block, whereas the flight this exists to stop leaves the
block on every call and ends every one of them \texttt{BUFFERED}. And the count must be
zeroed for reused photon slots, which are not value-initialised, or a new photon inherits
its predecessor's count.

The cap is checked inside the pusher loop, which bounds that loop, and again after it,
which catches the step where the photon left the block and the loop exited on its own
condition. Without the second check such a photon is shipped to the next block, possibly
through an MPI message, only to be retired on its first step there.

\paragraph{This is a biased estimator.} It discards photons, and not at random: it discards
exactly the longest-path ones. In a scattering-dominated test with $\tau = 3$ and roughly
ten scatterings per photon, a cap of $200$ steps still discarded $8.7\%$ of photons,
because free paths are long \emph{horizontally} even when the vertical depth is not ---
periodic directions have no bounding depth. It is a safety valve for genuinely unbounded
flights, not a general speed control, and wants to sit well above the longest physical free
flight. The default is $0$, meaning off, and $n_{\rm rem}$ reports how much was discarded.

\section{Resident photons}
\label{sec:resident}

A round pushes the photons currently resident, and how many that is sets both the working
set and the size the exchange buffers grow to. \texttt{loop\_max\_size} caps the photons
resident on one \emph{block}, so the memory it authorises is that figure multiplied by the
number of blocks a rank holds. On a mesh of a few large blocks this is invisible. On a mesh
of many small ones it is not: $10^{4}$ photons on each of $2601$ blocks is
$2.6\times10^{7}$ resident photons, several gigabytes, on blocks holding $256$ cells
apiece.

\texttt{max\_resident\_photons} bounds the product instead --- which is the quantity that
has to fit in memory --- by sharing a per-rank budget among that rank's blocks,
\begin{equation}
  \texttt{loop\_max\_size} \leftarrow
    \min\left(\texttt{loop\_max\_size},\;
      \max\left(1,\;
        \frac{\texttt{max\_resident\_photons}}{N^{\rm rank}_{\rm block}}\right)\right),
  \label{eq:budget}
\end{equation}
the per-block cap still applying on top, so a rank holding few blocks keeps exactly the
behaviour it had and existing results do not move. Neither figure changes any physics:
\texttt{loop\_max\_size} sets only how many photons a block emits per pass, the remainder
staying in \texttt{nphremain} for later ones. Measured on \texttt{mc\_isoth} with $512$
blocks and $10^{7}$ photons, where the $10^{4}$ cap does bind, the budget cut the peak
resident set from $2.06$ to $1.17$\,GB with the escaping fraction unchanged to $0.1\%$. The
reduction is reported at setup whenever it applies, so the choice is not silent.

\section{Parameters}

All in the \texttt{<montecarlo>} block.

\begin{center}
\begin{tabular}{lll}
\toprule
Name & Default & Meaning \\
\midrule
\texttt{rank\_exchange}   & on  & aggregate off-rank photons per rank, not per block link \\
\texttt{async\_term}      & on  & counter-based termination instead of a collective per round \\
\texttt{local\_max\_sweeps} & 1000 & consecutive same-rank transport sweeps per round \\
\texttt{capmove}          & 0 (off) & cap on steps in one free flight; see Sec.~\ref{sec:capmove} \\
\texttt{loop\_max\_size}  & 10000 & photons resident on one block; see Sec.~\ref{sec:resident} \\
\texttt{max\_resident\_photons} & $10^{6}$ & per-rank budget shared among that rank's blocks \\
\bottomrule
\end{tabular}
\end{center}

One further cap sits on the output side rather than in this block:
\texttt{<outputN>/max\_resident} (default $2\times10^{6}$) is how many escaped photons a
photon list holds before spilling to its file. Without it the list grows with the number of
photons that escape, which is unbounded in \texttt{nphot}, and is then copied whole in order
to be byte-swapped.

\texttt{async\_term} falls back to the synchronous loop whenever the rank exchange is
inactive, which includes every serial run. Setting it off is also how bit-reproducible
output is recovered when one is bisecting a change; see gap~2. \texttt{local\_max\_sweeps} bounds a same-rank
ping-pong; in the asynchronous loop it also ensures a return to the drain, so photons
arriving from other ranks do not wait behind an unbounded local sweep.

\section{Known gaps}

\begin{enumerate}
\item \textbf{Two transports.} Secs.~\ref{sec:sync} and \ref{sec:async} are independent
implementations. They now share the delivery walk of Sec.~\ref{sec:layout}, which is the
part that silently corrupts, but not the message structure. The synchronous path is worth
keeping only until the asynchronous one has enough mileage to retire it.

\item \textbf{Results are not reproducible.} Photons arrive in a different order, so each
block consumes its random stream differently. This holds between the two transports, and
it also holds for the asynchronous transport against itself: nothing pins the order in
which messages are taken, so two identical runs diverge. Measured on \texttt{mc\_isoth} at
four ranks with $4\times10^{5}$ photons, three repeats of one configuration gave
$n_{\rm esc} = 95898$, $95893$, $95898$, with the photon lists differing byte for byte ---
including one rank whose file was the same size both times but held different values, so
it is individual photons that move, not just how many escape. The same three repeats under
\texttt{async\_term = false} gave $95903$ every time, with files identical on all four
ranks.

Escaping fractions agree within Poisson noise, which is the most that can be asked of the
asynchronous path, and the conservation identity of Sec.~\ref{sec:conservation} is exact
either way. The practical consequence is for debugging rather than for physics: an MPI
regression cannot be found by diffing outputs unless the synchronous path is selected
first, and a run-to-run spread of a few tens of photons is easily mistaken for a real
change in behaviour. Any regression test over MPI must therefore either pin
\texttt{async\_term = false} or compare statistically.

\item \textbf{Load balance.} Blocks are distributed by cell count, and photons by emission,
but transport cost follows where photons travel and how much they scatter. Neither is
balanced by the current decomposition, and no improvement to the exchange addresses it.
Making the Athena++ load balancer aware of Monte Carlo work is the next task.

\item \textbf{Static peer set.} As in Sec.~\ref{sec:layout}, and a direct constraint on the
previous item: a load balancer that moves blocks mid-transport must rebuild the peer list
and the staged headers.
\end{enumerate}

\end{document}
